% comandi.tex — i comandi "\" dell'editor
%
% Nell'editor scrivi \ e l'inizio del nome: compare la lista, Invio (o Tab)
% sostituisce il comando con il blocco Quarto qui sotto. Il file e' tuo:
% aggiungi, togli, cambia. Si ricarica da solo a ogni apertura dell'editor.
%
%   \newcommand{\nome}{descrizione}
%   ...il testo che entra nel documento...
%   \endcommand
%
% \cursor segna dove finisce il cursore dopo l'inserimento.
% Le righe fuori da \newcommand ... \endcommand sono commenti.

% --- Struttura -----------------------------------------------------------------

\newcommand{\section}{Titolo di sezione (##)}
## \cursor
\endcommand

\newcommand{\subsection}{Sottosezione (###)}
### \cursor
\endcommand

\newcommand{\label}{Etichetta per i rimandi, da mettere dopo un titolo}
{#sec-\cursor}
\endcommand

\newcommand{\ref}{Rimando a un'etichetta: @sec-nome}
@sec-\cursor
\endcommand

\newcommand{\footnote}{Nota a piè di pagina}
^[\cursor]
\endcommand

% --- Testo ---------------------------------------------------------------------

\newcommand{\textbf}{Grassetto}
**\cursor**
\endcommand

\newcommand{\emph}{Corsivo}
*\cursor*
\endcommand

\newcommand{\texttt}{Codice in linea}
`\cursor`
\endcommand

\newcommand{\href}{Link}
[\cursor](https://)
\endcommand

\newcommand{\quote}{Citazione}
> \cursor
\endcommand

\newcommand{\itemize}{Elenco puntato}
- \cursor
-
\endcommand

\newcommand{\enumerate}{Elenco numerato}
1. \cursor
2.
\endcommand

% --- Matematica, figure, tabelle ---------------------------------------------------

\newcommand{\equation}{Equazione numerata}
$$
\cursor
$$ {#eq-nome}
\endcommand

\newcommand{\includegraphics}{Immagine con didascalia}
![\cursor](images/nome-file.png){fig-alt="" width=80%}
\endcommand

\newcommand{\tabular}{Tabella markdown}
| Modello | MAE | Baseline |
|---|---:|---:|
| \cursor |  |  |
\endcommand

\newcommand{\lstlisting}{Blocco di codice Python}
```python
\cursor
```
\endcommand

\newcommand{\embed}{Cella di un notebook in posts/_notebooks/}

\endcommand

% --- Blocchi del sito (le classi stanno in styles.css) --------------------------------

\newcommand{\hero}{Apertura grande, come in home}
::: {.hero}

::: {.hero-meta}
Milan, Italy · \cursor
:::

# Titolo {.no-anchor}

::: {.lede}
Due righe che dicono chi sei e cosa fai.
:::

:::
\endcommand

\newcommand{\sectionrule}{Titoletto di sezione con filetto, come in home}
::: {.section-rule}
\cursor
:::
\endcommand

\newcommand{\projectcard}{Scheda progetto}
::: {.project-card}
### [\cursor](projects/slug/index.qmd)

Una frase su cosa hai costruito e a quale domanda risponde.

::: {.stack}
Python · pandas
:::
:::
\endcommand

\newcommand{\stack}{Riga di tecnologie in monospazio}
::: {.stack}
\cursor
:::
\endcommand

\newcommand{\tag}{Etichetta piccola}
[\cursor]{.tag}
\endcommand

\newcommand{\morelink}{Link "vedi tutto →"}
[\cursor →](writing.qmd){.more-link}
\endcommand

\newcommand{\readinglist}{Lista di libri, come in Reading}
<ul class="reading-list">
<li><span class="status now">Reading</span><span><em>\cursor</em> <span class="book-author">— Autore</span></span></li>
<li><span class="status">Finished</span><span><em>Titolo</em> <span class="book-author">— Autore</span></span></li>
</ul>
\endcommand

% --- Impaginazione Quarto ------------------------------------------------------------

\newcommand{\callout}{Riquadro nota (note, tip, warning, important)}
::: {.callout-note}
## \cursor
:::
\endcommand

\newcommand{\columns}{Due colonne affiancate}
:::: {.columns}
::: {.column width="50%"}
\cursor
:::
::: {.column width="50%"}

:::
::::
\endcommand

\newcommand{\marginpar}{Nota a margine}
::: {.column-margin}
\cursor
:::
\endcommand

\newcommand{\widefigure}{Blocco piu' largo della colonna di testo}
::: {.column-page}
\cursor
:::
\endcommand

\newcommand{\newpage}{Filetto orizzontale}
---
\cursor
\endcommand
